
% SGA 3，第 VII 篇_B；英文重建本第一轮的中文译本。
% 范围： 节 4.3 开头 全文 依据文献局部页 76--77,
% 止于下节之前： 节 4.3.1.
% 合并阅读本页： 589--590. 重版印刷页： 578--579.
% 范围内原版边码： 545.
% 依据： Polo--Gille Exp7B-13oct24.pdf,
% SHA-256 D303690E0195654BBE7C3DAA5D89EA23E24CD709A3575217234AC19AE8836694.
% 参照了 Jacob Reinhold 的 CC BY 4.0 Markdown，作为对照来源。
% 本组件无图；所有数学内容均为原生 LaTeX。
\subsection*{4.3}
\phantomsection
\label{sga3:VIIB:4.3}
现在更详细地研究 \(\mathbf L\) 是 \(\mathbf O_k\) 上的平坦 \(p\)-Lie 代数时的形式 \(k\)-群 \(\mathscr G_p(\mathbf L)\)。
为此，先设 \(C\) 是特征 \(p\) 的环，\(L\) 是 \(C\) 上的 \(p\)-Lie 代数，其底模是以
\[
  (x_i)_{i\in I}.
\]
为基的自由模。令 \(\mathcal P\) 为非负整数族
\[
  n=(n_i)_{i\in I}
\]
的集合，满足
\[
  0\leq n_i<p
\]
且除可能有限多个指标外均有 \(n_i=0\)。赋予 \(I\) 一个全序，并令
\[
  i_1<i_2<\cdots<i_r
\]
为满足 \(n_i\ne0\) 的指标。置
\[
  |n|=n_{i_1}+\cdots+n_{i_r},
\]
以及
\[
  x^n
  =
  x_{i_1}^{\,n_{i_1}}
  x_{i_2}^{\,n_{i_2}}
  \cdots
  x_{i_r}^{\,n_{i_r}},
  \qquad
  n!
  =
  (n_{i_1}!)\cdots(n_{i_r}!).
\]
由 VII\textsubscript{A}，\SGARefLink{sga3:VIIA:5.3.3}{5.3.3}，单项式 \(x^n/n!\) 构成 \(U_p(L)\) 的一组基，且
\begin{equation}
  \Delta\left(\frac{x^n}{n!}\right)
  =
  \sum_{m\leq n}
  \frac{x^m}{m!}
  \otimes
  \frac{x^{\,n-m}}{(n-m)!}.
  \tag{*}
  \label{sga3:VIIB:4.3:restricted-pbw-comultiplication}
\end{equation}
求和遍历所有满足 \(m\leq n\) 的 \(m\in\mathcal P\)，即
\[
  m_i\leq n_i
  \qquad\text{对每个 }i\in I.
\]
公式\textup{(*)} 立即给出 \(U_p(L)\) 的自然滤过的描述\textup{（参见\ \SGARefLink{sga3:VIIB:1.3.6}{1.3.6}）}。置
\[
  U=U_p(L),
\]
令 \(U^+\) 为 \(L\) 生成的双边理想，并置
\[
  U_0=C\cdot1_U.
\]
与 \SGARefLink{sga3:VIIB:1.3.6}{1.3.6} 一样，定义 \(U\) 的 \(C\)-子余代数滤过：对 \(n\geq1\)，置
\[
  U_n
  =
  \left\{
    u\in U
    \,\middle|\,
    \Delta_U(u)-u\otimes1
    \in U_{n-1}\otimes U^+
  \right\}.
\]
于是公式\textup{(*)} 表明，\(U_n\) 是以满足
\[
  |m|\leq n.
\]
的单项式 \(x^m\) 为基的自由 \(C\)-模。
